%=====================================================================
%  IV. KESIMPULAN
%=====================================================================
\section{KESIMPULAN}

Berdasarkan penelitian yang telah dilakukan, diperoleh beberapa kesimpulan sebagai
berikut:

\begin{enumerate}
  \item Persamaan perturbasi medan skalar pada ruang-waktu Schwarzschild dapat
        diturunkan secara analitik dari persamaan Klein-Gordon dalam ruang-waktu
        lengkung. Sifat statis dan simetri bola latar belakang memungkinkan
        separasi variabel waktu dan sudut, sedangkan penskalaan fungsi radial
        $R(r) = \Psi(r)/r$ bersama koordinat \textit{tortoise} membawa persamaan
        radial ke dalam bentuk persamaan master yang menyerupai persamaan
        Schr\"{o}dinger tak gayut waktu, sehingga seluruh pengaruh geometri
        lubang hitam terangkum di dalam sebuah potensial efektif tunggal.

  \item Potensial efektif medan skalar tak bermassa berbentuk penghalang tunggal
        yang lenyap pada horizon peristiwa maupun pada jarak tak hingga.
        Ketinggian puncaknya meningkat secara signifikan terhadap bilangan
        momentum sudut akibat kontribusi suku sentrifugal, sedangkan letak
        puncaknya bergeser mendekati radius orbit foton $r = 3M$ untuk $\ell$ yang
        besar. Terhadap massa lubang hitam, potensial hanya bergantung pada
        kombinasi $r/M$, sehingga penambahan massa merendahkan dan melebarkan
        penghalang tanpa mengubah bentuknya.

  \item Kehadiran massa medan skalar mengubah perilaku asimtotik potensial efektif,
        yakni nilainya pada jarak jauh tidak lagi menuju nol melainkan menuju
        $\mu^{2}$. Akibatnya, penghalang potensial berdiri di atas dataran setinggi
        $\mu^{2}$ dan medan skalar bermassa tidak dapat merambat bebas menuju tak
        hingga. Pada nilai $\mu$ yang cukup besar, struktur penghalang tunggal
        bahkan lenyap dan potensial berubah menjadi fungsi yang naik secara
        monoton dari horizon peristiwa menuju dataran asimtotiknya.
\end{enumerate}
